\begin{abstract}
A formula-first co-design is proposed for a non-ideal asynchronous Buck
converter fed at \SI{48}{V} and commanded over \SIrange{1}{46}{V}. The
search space is the set of Bryson excursions of a linear--quadratic--integral
(LQI) cost; each candidate is solved from the Riccati equation and realized
exactly as a two-degree-of-freedom PIDF that measures the output voltage only.
We prove that the derivative filter, with its pole on the ESR zero,
reconstructs the capacitor voltage along every trajectory, so the PIDF and the full-state LQI coincide under saturation,
anti-windup and discontinuous conduction and differ only through an explicit
output-current feedback, an output-network mismatch and sampling; a
triangle-hold discretization keeps the sampled PIDF within \EqSampFOH{} of the
sampled LQI on a \SI{45}{V} step, against \EqSampZOH{} with a zero-order hold.
Six equal-budget optimizers with Gaussian-process refinement search the weights
under 33 inequalities in engineering units, three measured on a switched model
because averaged-feasible designs hold coexisting switching limit cycles. Over
30 matched seeds all optimizers reach the same optimum, \CMOgainPct\,\% below
the Bryson design. A direct gain-space search reaches the same LQ-optimal controller, but
\IOpct\,\% of its feasible designs are not LQ-optimal. On a switched replica of the reference Simulink model, the retained
controller removes the overshoot of classical designs and stays within the
device current rating; its only miss is \SwEss{}\,mV of steady error at
\SI{1}{V} against \SI{25}{mV}, which its full-state twin meets; standard PI, PID, PIDF, LQR, LQG and LQI designs exceed the specification by factors of 6 to 30.
\end{abstract}
